\documentclass[10pt]{beamer}
\usetheme{metropolis}
\usepackage{graphicx}
\usepackage{listings}
\usepackage{booktabs}
\usepackage{xcolor}
\usepackage{hyperref}
\definecolor{codebg}{HTML}{F5F5F5}
\definecolor{codegreen}{HTML}{2E7D32}
\definecolor{codegray}{HTML}{757575}
\definecolor{codeblue}{HTML}{1565C0}
\lstdefinestyle{rstyle}{
  language=R, backgroundcolor=\color{codebg}, basicstyle=\ttfamily\scriptsize,
  keywordstyle=\color{codeblue}\bfseries, stringstyle=\color{codegreen},
  commentstyle=\color{codegray}\itshape, breaklines=true, frame=single,
  rulecolor=\color{codegray!30}, showstringspaces=false, tabsize=2,
  morekeywords={library,ggplot,aes,geom_histogram,geom_density,geom_freqpoly,
    geom_density_ridges,stat_ecdf,scale_fill_manual,scale_color_manual,
    theme_minimal,theme,element_text,labs,after_stat,facet_wrap}
}
\lstdefinestyle{pythonstyle}{
  language=Python, backgroundcolor=\color{codebg}, basicstyle=\ttfamily\scriptsize,
  keywordstyle=\color{codeblue}\bfseries, stringstyle=\color{codegreen},
  commentstyle=\color{codegray}\itshape, breaklines=true, frame=single,
  rulecolor=\color{codegray!30}, showstringspaces=false, tabsize=4,
  morekeywords={import,from,as,pd,plt,sns,np,fig,ax,hist,plot,fill_between,
    gaussian_kde,histplot,kdeplot,ecdfplot,FacetGrid}
}
\title{Distributions: Histograms, Density \& Ridgeline}
\subtitle{Workshop 3 --- Module 1}
\author{Prof.Asoc.\ Endri Raco}
\institute{Data Visualization for Data Scientists \\ UNYT Tirana}
\date{2025/26}
\begin{document}

\begin{frame}
  \titlepage
\end{frame}

\begin{frame}{Module Roadmap}
  \begin{enumerate}
    \item Histograms: bin width selection and its consequences
    \item Histogram variants: frequency, density, cumulative
    \item Kernel Density Estimation: theory and bandwidth
    \item Histogram + density overlay
    \item Grouped densities and ridgeline plots
    \item Distribution chart selection guide
  \end{enumerate}
  \begin{exampleblock}{Learning Outcome}
    You will choose the correct distribution chart for any dataset,
    understand why bin width and bandwidth matter, and produce
    ridgeline plots for multi-group comparisons.
  \end{exampleblock}
\end{frame}

\section{Histograms}

\begin{frame}{Bin Width Matters}
  \begin{center}
    \includegraphics[width=0.95\textwidth]{figures_w03m01/bin_width}
  \end{center}
  \begin{alertblock}{Pro-Tip}
    Too few bins hides structure; too many creates noise. Sturges' rule
    ($k = 1 + \log_2 n$) and Freedman-Diaconis ($h = 2 \cdot IQR \cdot n^{-1/3}$)
    provide starting points. Always try 2--3 values and inspect.
  \end{alertblock}
\end{frame}

\begin{frame}{Three Histogram Variants}
  \begin{center}
    \includegraphics[width=0.92\textwidth]{figures_w03m01/hist_variants}
  \end{center}
  \begin{exampleblock}{Data Insight}
    \textbf{Frequency}: raw counts, intuitive but scale-dependent.
    \textbf{Density}: area sums to 1, comparable across sample sizes.
    \textbf{Cumulative}: shows percentile structure (``50\% of apps
    are rated below 4.2'').
  \end{exampleblock}
\end{frame}

\begin{frame}[fragile]{Histograms in R vs Python}
  \begin{columns}[T]
    \begin{column}{0.48\textwidth}
      \textbf{R (ggplot2)}
\begin{lstlisting}[style=rstyle]
# Frequency histogram
ggplot(apps, aes(x = Rating)) +
  geom_histogram(bins = 30,
    fill = "#1565C0",
    color = "white") +
  theme_minimal()

# Density histogram (area = 1)
ggplot(apps, aes(x = Rating)) +
  geom_histogram(
    aes(y = after_stat(density)),
    bins = 30, fill = "#1565C0",
    color = "white") +
  theme_minimal()

# Cumulative (ECDF)
ggplot(apps, aes(x = Rating)) +
  stat_ecdf(color = "#E65100",
            linewidth = 1) +
  theme_minimal()
\end{lstlisting}
    \end{column}
    \begin{column}{0.48\textwidth}
      \textbf{Python}
\begin{lstlisting}[style=pythonstyle]
# seaborn (recommended)
sns.histplot(apps, x="Rating",
    bins=30, kde=False)

# Density histogram
sns.histplot(apps, x="Rating",
    bins=30, stat="density")

# ECDF
sns.ecdfplot(apps, x="Rating")

# matplotlib (manual)
ax.hist(ratings, bins=30,
    color="#1565C0",
    edgecolor="white")

# Density
ax.hist(ratings, bins=30,
    density=True)

# Cumulative
ax.hist(ratings, bins=50,
    cumulative=True, density=True)
\end{lstlisting}
    \end{column}
  \end{columns}
\end{frame}

\section{Kernel Density Estimation}

\begin{frame}{KDE: How It Works}
  \begin{center}
    \includegraphics[width=0.92\textwidth]{figures_w03m01/kde_concept}
  \end{center}
\end{frame}

\begin{frame}{Bandwidth Controls Smoothness}
  \begin{center}
    \includegraphics[width=0.92\textwidth]{figures_w03m01/bandwidth}
  \end{center}
  \begin{alertblock}{Pro-Tip}
    Bandwidth is to KDE what bin width is to histograms. Too small
    (undersmoothed): noise. Too large (oversmoothed): hides modes.
    Default bandwidth selectors (Silverman's rule) work well for
    unimodal data but may oversmooth multimodal distributions.
  \end{alertblock}
\end{frame}

\begin{frame}[fragile]{KDE in R vs Python}
  \begin{columns}[T]
    \begin{column}{0.48\textwidth}
      \textbf{R (ggplot2)}
\begin{lstlisting}[style=rstyle]
# Density curve
ggplot(apps, aes(x = Rating)) +
  geom_density(fill = "#1565C0",
    alpha = 0.3, linewidth = 1) +
  theme_minimal()

# Adjust bandwidth
ggplot(apps, aes(x = Rating)) +
  geom_density(
    adjust = 0.5,  # half default bw
    fill = "#1565C0",
    alpha = 0.3) +
  theme_minimal()

# adjust > 1 = smoother
# adjust < 1 = rougher
\end{lstlisting}
    \end{column}
    \begin{column}{0.48\textwidth}
      \textbf{Python}
\begin{lstlisting}[style=pythonstyle]
# seaborn
sns.kdeplot(apps, x="Rating",
    fill=True, alpha=0.3,
    color="#1565C0")

# Adjust bandwidth
sns.kdeplot(apps, x="Rating",
    bw_adjust=0.5,  # half default
    fill=True, alpha=0.3)

# scipy (manual)
from scipy.stats import gaussian_kde
kde = gaussian_kde(ratings,
    bw_method=0.3)
xs = np.linspace(1, 5, 200)
ax.plot(xs, kde(xs))
ax.fill_between(xs, kde(xs),
    alpha=0.15)
\end{lstlisting}
    \end{column}
  \end{columns}
\end{frame}

\section{Histogram + Density Overlay}

\begin{frame}{The Best of Both: Histogram + KDE + Reference Lines}
  \begin{center}
    \includegraphics[width=0.78\textwidth]{figures_w03m01/hist_density_overlay}
  \end{center}
  \begin{exampleblock}{Data Insight}
    Overlaying KDE on a density-scaled histogram shows both the binned
    structure and the smooth shape. Adding mean and median lines reveals
    skewness: when mean $<$ median, the distribution is left-skewed.
  \end{exampleblock}
\end{frame}

\begin{frame}[fragile]{Overlay in R vs Python}
  \begin{columns}[T]
    \begin{column}{0.48\textwidth}
\begin{lstlisting}[style=rstyle]
ggplot(apps, aes(x = Rating)) +
  geom_histogram(
    aes(y = after_stat(density)),
    bins = 35, fill = "#1565C0",
    color = "white", alpha = 0.5) +
  geom_density(color = "#E53935",
    linewidth = 1.2) +
  geom_vline(
    xintercept = mean(apps$Rating,
      na.rm = TRUE),
    linetype = "dashed",
    color = "#2E7D32", linewidth = 1) +
  theme_minimal() +
  labs(title = "Histogram + KDE",
       x = "Rating", y = "Density")
\end{lstlisting}
    \end{column}
    \begin{column}{0.48\textwidth}
\begin{lstlisting}[style=pythonstyle]
# seaborn: one line!
sns.histplot(apps, x="Rating",
    bins=35, kde=True,
    stat="density",
    color="#1565C0", alpha=0.5)

# Or matplotlib manual:
ax.hist(ratings, bins=35,
    density=True, color="#1565C0",
    edgecolor="white", alpha=0.5)
kde = gaussian_kde(ratings)
xs = np.linspace(1, 5, 200)
ax.plot(xs, kde(xs),
    color="#E53935", lw=2)
ax.axvline(x=np.mean(ratings),
    color="#2E7D32", lw=1.5,
    linestyle="--")
\end{lstlisting}
    \end{column}
  \end{columns}
\end{frame}

\section{Grouped Density \& Ridgeline}

\begin{frame}{Grouped Density: 2--3 Groups}
  \begin{center}
    \includegraphics[width=0.78\textwidth]{figures_w03m01/grouped_density}
  \end{center}
\end{frame}

\begin{frame}[fragile]{Grouped Density in Code}
  \begin{columns}[T]
    \begin{column}{0.48\textwidth}
\begin{lstlisting}[style=rstyle]
ggplot(apps, aes(x = Rating,
    fill = Type, color = Type)) +
  geom_density(alpha = 0.2,
    linewidth = 1) +
  scale_fill_manual(values =
    c(Free = "#1565C0",
      Paid = "#E53935")) +
  scale_color_manual(values =
    c(Free = "#1565C0",
      Paid = "#E53935")) +
  theme_minimal() +
  labs(title = "Free vs Paid")
\end{lstlisting}
    \end{column}
    \begin{column}{0.48\textwidth}
\begin{lstlisting}[style=pythonstyle]
# seaborn
sns.kdeplot(data=apps, x="Rating",
    hue="Type", fill=True,
    alpha=0.2, linewidth=1.5,
    palette={"Free":"#1565C0",
             "Paid":"#E53935"})

# matplotlib manual
for t, c in [("Free","#1565C0"),
             ("Paid","#E53935")]:
    sub = apps.query("Type==@t")
    kde = gaussian_kde(
        sub["Rating"].dropna())
    ax.fill_between(xs, kde(xs),
        alpha=0.15, color=c)
    ax.plot(xs, kde(xs),
        color=c, lw=1.5, label=t)
\end{lstlisting}
    \end{column}
  \end{columns}
\end{frame}

\begin{frame}{Ridgeline Plot: 4+ Groups}
  \begin{center}
    \includegraphics[width=0.78\textwidth]{figures_w03m01/ridgeline}
  \end{center}
  \begin{alertblock}{Pro-Tip}
    Ridgeline plots (joy plots) stack densities vertically, using
    slight overlap for compactness. In R: \texttt{ggridges::geom\_density\_ridges()}.
    In Python: manual offset loop or \texttt{joypy} package.
    Use when grouped density would produce too much overlap ($>$3 groups).
  \end{alertblock}
\end{frame}

\begin{frame}[fragile]{Ridgeline in R}
\begin{lstlisting}[style=rstyle]
library(ggridges)

ggplot(apps, aes(x = Rating,
                  y = `Content Rating`,
                  fill = `Content Rating`)) +
  geom_density_ridges(
    alpha = 0.4,
    scale = 1.5,          # overlap amount
    rel_min_height = 0.01 # trim tails
  ) +
  scale_fill_brewer(palette = "Set2") +
  theme_minimal() +
  theme(legend.position = "none") +
  labs(title = "Rating by Content Rating",
       x = "Rating", y = NULL)
\end{lstlisting}
\end{frame}

\section{Distribution Selection}

\begin{frame}{Chart Selection Guide}
  \begin{center}
    \includegraphics[width=0.92\textwidth]{figures_w03m01/decision_tree}
  \end{center}
\end{frame}

\section{Complete Examples}

\begin{frame}[fragile]{R: Publication-Quality Distribution Chart}
  \begin{columns}[T]
    \begin{column}{0.52\textwidth}
\begin{lstlisting}[style=rstyle]
ggplot(apps |>
    filter(!is.na(Rating)),
  aes(x = Rating)) +
  geom_histogram(
    aes(y = after_stat(density)),
    bins = 35, fill = "#1565C0",
    color = "white", alpha = 0.5) +
  geom_density(color = "#E53935",
    linewidth = 1.2) +
  geom_vline(
    xintercept = mean(
      apps$Rating, na.rm = TRUE),
    linetype = "dashed",
    color = "#2E7D32") +
  annotate("text",
    x = 3.5, y = 1.2,
    label = "Left-skewed:\nhigh-rating mode",
    fontface = "italic",
    color = "#555", size = 3) +
  theme_minimal() +
  labs(title = "App Ratings are
       Left-Skewed",
    subtitle = "Google Play Store,
                n = 9,366",
    x = "Rating", y = "Density")
\end{lstlisting}
    \end{column}
    \begin{column}{0.45\textwidth}
      \includegraphics[width=\textwidth]{figures_w03m01/r_dist}
    \end{column}
  \end{columns}
\end{frame}

\begin{frame}[fragile]{Python: Same Chart}
  \begin{columns}[T]
    \begin{column}{0.52\textwidth}
\begin{lstlisting}[style=pythonstyle]
fig, ax = plt.subplots(figsize=(6,5))

ax.hist(ratings, bins=35,
    density=True, color="#2E7D32",
    edgecolor="white", lw=0.3,
    alpha=0.5)

kde = gaussian_kde(ratings)
xs = np.linspace(0.5, 5.5, 200)
ax.plot(xs, kde(xs),
    color="#E53935", lw=2)

ax.axvline(x=np.mean(ratings),
    color="#2E7D32", lw=1.5,
    linestyle="--")

ax.set_title("App Ratings are
    Left-Skewed",
    fontweight="bold")
ax.set_xlabel("Rating")
ax.set_ylabel("Density")
ax.spines["top"].set_visible(False)
ax.spines["right"].set_visible(False)
\end{lstlisting}
    \end{column}
    \begin{column}{0.45\textwidth}
      \includegraphics[width=\textwidth]{figures_w03m01/py_dist}
    \end{column}
  \end{columns}
\end{frame}

\section{Synthesis}

\begin{frame}{Key Takeaways}
  \begin{enumerate}
    \item \textbf{Bin width} determines histogram structure: always try
          2--3 values. Use Freedman-Diaconis as a starting point.
    \item \textbf{KDE} smooths the histogram; \textbf{bandwidth} controls
          smoothness. Default rules work for unimodal data.
    \item \textbf{Overlay} histogram + KDE + reference lines for the most
          informative single-variable display.
    \item \textbf{Grouped density} works for $\leq$3 groups;
          \textbf{ridgeline} for 4+ groups.
    \item In R: \texttt{geom\_histogram}, \texttt{geom\_density},
          \texttt{ggridges::geom\_density\_ridges}.
          In Python: \texttt{sns.histplot(kde=True)},
          \texttt{sns.kdeplot}, \texttt{joypy}/manual offset.
  \end{enumerate}
\end{frame}

\begin{frame}{Essential Readings}
  \begin{itemize}
    \item Wickham, H. (2016). \emph{ggplot2}, Chapter 5 (Statistical Summaries).
    \item Wilke, C.~O. (2019). \emph{Fundamentals of Data Visualization},
          Chapters 7--9 (Histograms, Density, Multiple Distributions).
    \item Silverman, B.~W. (1986). \emph{Density Estimation for Statistics
          and Data Analysis}. Chapman \& Hall.
  \end{itemize}
  \vspace{6pt}
  \textbf{Next Module}: Distributions: Boxplots, Violins \& Beeswarm (W03-M02)
\end{frame}

\end{document}
